\begin{frame}{Main Theorem}
    \begin{theorem}
        Let $\Lambda$ be a unimodal logic with $\Lambda \in \{T,D\}$. Then the following holds:
        \[
           \Lambda \times_n^+ \Lambda = \Lambda \otimes \Lambda \otimes \Lambda + \Box p \to \Box_1\Box_2p \land \Box_2\Box_1p 
        \]
    \end{theorem}
    In the following we will sketch the proof for $\Lambda = T$. \pause 
	Soundness ("$\supseteq$") is straightforward. 

        \vspace{1.5cm}
         We denote $T \otimes T \otimes T + \Box p \to \Box_1\Box_2p \land \Box_2\Box_1p$ as $T_3m$

\end{frame}
    
\begin{frame}{Proof steps}
    \begin{itemize}
        \item Idea: create an infinite branching and infinite depth Kripke tree with three relations (called $T_{\omega,\omega,\omega}^{r,m}$) with $Log(T_{\omega,\omega,\omega}^{r,m}) = T_3m$ 
		%\pause
        %\item Note: For any Kripke frame $F = (W, R_1,...R_n)$, we can construct a neighborhood frame $F'$, i.e.
      %  for any $w \in W$: $\tau_i(w) = \{U \mid R_i(w) \subseteq U \subseteq W\}$ where $Log(F) = Log(F')$ \pause
      %  \item construct a neighbourhood product $N_{\omega}^r \times_n^+ N_{\omega}^r$ s.t. there is surjective bounded morphism from $N_{\omega}^r \times_n^+ N_{\omega}^r$ to $N(T_{\omega,\omega,\omega}^{r,m})$
    \end{itemize}

	\vspace{2cm}
\[ 
	{\transparent{0.3} T \otimes T \otimes T + (\Box p \to \Box_1\Box_2 p \land \Box_2\Box_1 p) = T_3m }
 \]


    

	

\end{frame}

\begin{frame}{Proof steps}
	\begin{itemize}
	\setlength{\itemsep}{1em}

	\item Idea: create an infinite branching and depth Kripke tree with three relations (called $T_{\omega,\omega,\omega}^{r,m}$) with $Log(T_{\omega,\omega,\omega}^{r,m}) = T_3m$
    \item construct a neighbourhood product $N_{\omega}^r \times_n^+ N_{\omega}^r$ with $Log(N_\omega^r) = T$ s.t. there is surjective bounded morphism from $N_{\omega}^r \times_n^+ N_{\omega}^r$ to $N(T_{\omega,\omega,\omega}^{r,m})$

			\end{itemize}
\vspace{1.2cm}
\[ 
	{\transparent{0.3} T \otimes T \otimes T + (\Box p \to \Box_1\Box_2 p \land \Box_2\Box_1 p) = T_3m }
 \]
	
\end{frame}

\begin{frame}{Proof steps}
	\begin{itemize}
		\setlength{\itemsep}{1em}
	    \item Idea: create an infinite branching and depth Kripke tree with three relations (called $T_{\omega,\omega,\omega}^{r,m}$) with $Log(T_{\omega,\omega,\omega}^{r,m}) = T_3m$ 
	    \item construct a neighbourhood product $N_{\omega}^r \times_n^+ N_{\omega}^r$ 
		with $Log(N_\omega^r) = T$ s.t. there is surjective bounded morphism from $N_{\omega}^r \times_n^+ N_{\omega}^r$ to $N(T_{\omega,\omega,\omega}^{r,m})$
        \item Note: For any Kripke frame $F = (W, R_1,...R_n)$, we can construct a neighborhood frame $N(F)$, i.e.
       for any $w \in W$: $\tau_i(w) = \{U \mid R_i(w) \subseteq U \subseteq W\}$ where $Log(F) = Log(N(F))$ 

		\end{itemize}

		\[ 
	{\transparent{0.3} T \otimes T \otimes T + (\Box p \to \Box_1\Box_2 p \land \Box_2\Box_1 p) = T_3m }
 \]
\end{frame}












\begin{frame}{Log($T_{\omega,\omega,\omega}^{r,m}) = T_3m$}
        \begin{figure}[t]
	\centering
	\begin{tikzpicture}[scale=0.9]
		%koren
		\node (K) at (0,0) [circle,fill=black,inner sep=1pt] {};
		%perviy uroven
		\node (A) at (-5,-2) [circle,fill=black,inner sep=1pt] {};
		\node (B) at (-2,-2) [circle,fill=black,inner sep=1pt] {};
		\node (C) at (-1,-2) [circle,fill=black,inner sep=1pt] {};
		\node (K2) at (0,-2) [circle,fill=black,inner sep=1pt] {};
		\node (D) at (1,-2) [circle,fill=black,inner sep=1pt] {};
		\node (E) at (2,-2) [circle,fill=black,inner sep=1pt] {};
		\node (F) at (5,-2) [circle,fill=black,inner sep=1pt] {};
		%zapolnaem promejutki melkimi tockami
        \draw[line width=0.8pt, dash pattern=on 0pt off 6pt, line cap=round] (-4.7,-2) -- (-2.3,-2);
        \draw[line width=0.8pt, dash pattern=on 0pt off 6pt, line cap=round] (-.7,-2) -- (.7,-2);
        \draw[line width=0.8pt, dash pattern=on 0pt off 6pt, line cap=round] (2.3,-2) -- (4.7,-2);
        %risuem linii mejdu kornem i pervim urovnem
        \draw[->, >=stealth, wred, thick](K) -- node[midway, left = 9pt] {\footnotesize $R_1$} (A);
        \draw[->, >=stealth, wred, thick] (K) -- (B);
        \draw[->, >=stealth, mblue, thick](K) -- node[midway, right = 0.1pt] {\footnotesize $R$} (C);
        \draw[->, >=stealth, mblue, thick] (K) -- (D);
        \draw[->, >=stealth, sgreen, thick] (K) -- (E);
        \draw[->, >=stealth, sgreen, thick](K) -- node[midway, right = 9pt] {\footnotesize $R_2$} (F);
        %versini dla 3qo urovna
        \node (A1) at (-7,-3.5) [circle,fill=black,inner sep=1pt] {};
        \node (B1) at (-6,-3.5) [circle,fill=black,inner sep=1pt] {};
        \node (C1) at (-5.5,-3.5) [circle,fill=black,inner sep=1pt] {};
        \node (D1) at (-4.5,-3.5) [circle,fill=black,inner sep=1pt] {};
        \node (E1) at (-4,-3.5) [circle,fill=black,inner sep=1pt] {};
        \node (F1) at (-3,-3.5) [circle,fill=black,inner sep=1pt] {};
		%risuem linii mejdu kornem i pervim urovnem
		\draw[->, >=stealth, wred, thick] (A) -- (A1);       
		\draw[->, >=stealth, wred, thick] (A) -- (B1);
		\draw[->, >=stealth, mblue, thick] (A) -- (C1);
		\draw[->, >=stealth, mblue, thick] (A) -- (D1);
		\draw[->, >=stealth, sgreen, thick] (A) -- (E1);
		\draw[->, >=stealth, sgreen, thick] (A) -- (F1);
		%zapolnaem promejutki melkimi tockami
		\draw[line width=0.8pt, dash pattern=on 0pt off 6pt, line cap=round] (-6.8,-3.5) -- (-6.1,-3.5);
		\draw[line width=0.8pt, dash pattern=on 0pt off 6pt, line cap=round] (-5.3,-3.5) -- (-4.3,-3.5);
		\draw[line width=0.8pt, dash pattern=on 0pt off 6pt, line cap=round] (-3.8,-3.5) -- (-3.1,-3.5);
		%novie linii zascet svoystv nasey logiki
		\draw[->, >=stealth, mblue, thick] (K) -- (E1);
		\draw[->, >=stealth, mblue, thick] (K) -- (F1);
		%versini dla 3qo urovna 2ya cast
		\node (A2) at (-2,-3.5) [circle,fill=black,inner sep=1pt] {};
		\node (B2) at (-1,-3.5) [circle,fill=black,inner sep=1pt] {};
		\node (C2) at (-.5,-3.5) [circle,fill=black,inner sep=1pt] {};
		\node (D2) at (.5,-3.5) [circle,fill=black,inner sep=1pt] {};
		\node (E2) at (1,-3.5) [circle,fill=black,inner sep=1pt] {};
		\node (F2) at (2,-3.5) [circle,fill=black,inner sep=1pt] {};
		%risuem linii mejdu kornem i pervim urovnem
		\draw[->, >=stealth, wred, thick] (K2) -- (A2);       
		\draw[->, >=stealth, wred, thick] (K2) -- (B2);
		\draw[->, >=stealth, mblue, thick] (K2) -- (C2);
		\draw[->, >=stealth, mblue, thick] (K2) -- (D2);
		\draw[->, >=stealth, sgreen, thick] (K2) -- (E2);
		\draw[->, >=stealth, sgreen, thick] (K2) -- (F2);
		%zapolnaem promejutki melkimi tockami
		\draw[line width=0.8pt, dash pattern=on 0pt off 6pt, line cap=round] (-1.8,-3.5) -- (-1.1,-3.5);
		\draw[line width=0.8pt, dash pattern=on 0pt off 6pt, line cap=round] (-.3,-3.5) -- (0.7,-3.5);
		\draw[line width=0.8pt, dash pattern=on 0pt off 6pt, line cap=round] (1.2,-3.5) -- (1.9,-3.5);
		%versini dla 3qo urovna 3ya cast
		\node (A3) at (3,-3.5) [circle,fill=black,inner sep=1pt] {};
		\node (B3) at (4,-3.5) [circle,fill=black,inner sep=1pt] {};
		\node (C3) at (4.5,-3.5) [circle,fill=black,inner sep=1pt] {};
		\node (D3) at (5.5,-3.5) [circle,fill=black,inner sep=1pt] {};
		\node (E3) at (6,-3.5) [circle,fill=black,inner sep=1pt] {};
		\node (F3) at (7,-3.5) [circle,fill=black,inner sep=1pt] {};
		%risuem linii mejdu kornem i pervim urovnem
		\draw[->, >=stealth, wred, thick] (F) -- (A3);       
		\draw[->, >=stealth, wred, thick] (F) -- (B3);
		\draw[->, >=stealth, mblue, thick] (F) -- (C3);
		\draw[->, >=stealth, mblue, thick] (F) -- (D3);
		\draw[->, >=stealth, sgreen, thick] (F) -- (E3);
		\draw[->, >=stealth, sgreen, thick] (F) -- (F3);
		%zapolnaem promejutki melkimi tockami
		\draw[line width=0.8pt, dash pattern=on 0pt off 6pt, line cap=round] (3.2,-3.5) -- (3.9,-3.5);
		\draw[line width=0.8pt, dash pattern=on 0pt off 6pt, line cap=round] (4.7,-3.5) -- (5.6,-3.5);
		\draw[line width=0.8pt, dash pattern=on 0pt off 6pt, line cap=round] (6.2,-3.5) -- (6.9,-3.5);
		%novie linii zascet svoystv nasey logiki
		\draw[->, >=stealth, mblue, thick] (K) -- (A3);
		\draw[->, >=stealth, mblue, thick] (K) -- (B3);
		
		\foreach \x in {-6.6,-3,0, 3, 6.5}
		{
			\foreach \y in {-3.7, -4, -4.3}
			{
				\fill (\x,\y) circle (0.75pt);
			}
		}
		
		
	\end{tikzpicture}
\end{figure}
\[ 
	{\transparent{0.3} T \otimes T \otimes T + \Box p \to \Box_1\Box_2 p \land \Box_2\Box_1 p = T_3m }
 \]
\pause

\begin{itemize}
    \item Note: $\Box p \to \Box_1 p \land \Box_2p$ is derivable
    \item Log($T_{\omega,\omega,\omega}^{r,m}) \subseteq T_3m$ by standard techniques
\end{itemize}
\end{frame}

\begin{frame}{Definition of the neighborhood frame $N_w^r$}
    \begin{definition}
        Let $T_{\omega}^{r} = (\mathbb{N}^*, R)$ be the infinite branching  and infinite depth tree. Then
		\[
		X = \{a_0a_1a_2... \in (\mathbb{N} \cup \star) \mid \exists N \in \mathbb{N} \mbox{ }\forall k\geq N: a_k = \star\}
		\]
		is the set of pseudo infinite sequences. \pause 
		Let $\alpha =a_0a_1a_2... \in X$. Then 
		\[
		U_k(\alpha) = \{\beta \in X \mid {\mbox{\emph{forg}}}(\alpha) R \mbox{}\mbox{\emph{forg}}(\beta) \mbox{ and } \alpha \mid_m = \beta \mid_m,	 \mbox{with } m = max(k,st(\alpha))\} 
		\]
		where \emph{forg}($\alpha$) deletes all $\star$ and $st(\alpha) = min\{N \mid \forall k \geq N: a_k = \star\}$
		%denote the set of pseudo infinite sequences
    \end{definition}
\end{frame}

\begin{frame}{Example of $U_k(\alpha)$}
	\begin{center}
		

	\begin{tikzpicture}[scale=2.5]
		\node (K) at (0,0) [circle,fill=black,inner sep=1pt] {};
		\node (A) at (-0.5,-1) [circle,fill=black,inner sep=1pt] {};
		\node (B) at (-1,-1) [circle,fill=black,inner sep=1pt] {};
		\node (C) at (0,-1) [circle,fill=black,inner sep=1pt] {};
		\node (D) at (0.5,-1) [circle,fill=black,inner sep=1pt] {};
		\node (E) at (1,-1) [circle,fill=black,inner sep=1pt] {};

		\draw[->, >=stealth, wred, thick] (K) -- (A);
		\draw[->, >=stealth, wred, thick] (K) -- (B);
		\draw[->, >=stealth, wred, thick] (K) -- (C);
		\draw[->, >=stealth, wred, thick] (K) -- (D);
		\draw[->, >=stealth, wred, thick] (K) -- (E);

		\node  at (-0.2,0) [circle,fill=black,inner sep=0.2pt] {};
		\node  at (-0.3,0) [circle,fill=black,inner sep=0.2pt] {};
		\node  at (-0.4,0) [circle,fill=black,inner sep=0.2pt] {};

		\node  at (0,0.1) [circle,fill=black,inner sep=0.2pt] {};
		\node  at (0,0.2) [circle,fill=black,inner sep=0.2pt] {};
		\node  at (0,0.3) [circle,fill=black,inner sep=0.2pt] {};

		\node  at (0,-1.1) [circle,fill=black,inner sep=0.2pt] {};
		\node  at (0,-1.2) [circle,fill=black,inner sep=0.2pt] {};
		\node  at (0,-1.3) [circle,fill=black,inner sep=0.2pt] {};

		\node  at (1.1,-1) [circle,fill=black,inner sep=0.2pt] {};
		\node  at (1.2,-1) [circle,fill=black,inner sep=0.2pt] {};
		\node  at (1.3,-1) [circle,fill=black,inner sep=0.2pt] {};

		\node  at (-1.1,-1) [circle,fill=black,inner sep=0.2pt] {};
		\node  at (-1.2,-1) [circle,fill=black,inner sep=0.2pt] {};
		\node  at (-1.3,-1) [circle,fill=black,inner sep=0.2pt] {};

		% label above the dot
		\node at (0.2,0) {$1234$};
		\node at (0.7,-1.1) {$12344$}; 


	\end{tikzpicture}
	\end{center}
	

	Let $\alpha = \star \star 12 \star 3 \star 4 \star \star^\omega \in X$. Then $st(\alpha) = 9$. Let's consider $U_9(\alpha)$.
	Then 
	\[
		\beta = \star \star 12 \star 3 \star 4 \star 4 \star \star^\omega \in U_9(\alpha)
	\]
\end{frame}

\begin{frame}{Definition of the neighborhood frame $N_w^r$}
    \begin{definition}
        Let $T_{\omega}^{r} = (\mathbb{N}^*, R)$ be the infinite branching depth tree. Then
		\[
		X = \{a_0a_1a_2... \in (\mathbb{N} \cup \star)^\omega \mid \exists N \in \mathbb{N} \mbox{ }\forall k\geq N: a_k = \star\}
		\]
		is the set of pseudo infinite sequences. 
		Let $\alpha =a_0a_1a_2... \in X$. Then 
		\[
		U_k(\alpha) = \{\beta \in X \mid {\mbox{\emph{forg}}}(\alpha) R \mbox{}\mbox{\emph{forg}}(\beta) \mbox{ and } \alpha \mid_m = \beta \mid_m,	 \mbox{with } m = max(k,st(\alpha))\} 
		\] 
		where \emph{forg}($\alpha$) deletes all $\star$ and $st(\alpha) = min\{N \mid \forall k \geq N: a_k = \star\}$ \newline \newline \pause
		\textcolor{red}{We define $N_\omega^r = (X,\tau)$ where $\tau: X \to 2^{2^X}$ with 
			\[
				\tau(\alpha) = \{U \subseteq X \mid \exists k \in \mathbb{N}: U_k(\alpha) \subseteq U\}
			\]}
			Note: $\tau$ is a filter and there is no smallest neighborhood
		
    \end{definition}
	%Note: $\tau$ is a filter and there is no smallest neighborhood
\end{frame}

\begin{frame}[t]{Surjective bounded morphism}
	
	Let $N_\omega^r \times_n^+ N_\omega^r = (X \times X, \tau_1, \tau_2, \tau)$ and $N(T_{\omega,\omega,\omega}^{r,m}) = ((\mathbb{N} \times \{0,1,2\})^*, \sigma_1, \sigma_2, \sigma)$\pause
	\begin{theorem}
	 There is a surjective bounded morphism $g: N_\omega^r \times_n^+ N_\omega^r \to N(T_{\omega,\omega,\omega}^{r,m})$. 
	\end{theorem}
%	Note: $Log(N_w^r) = T$ and hence $Log(N_w^r \times_n^+ N_w^r) \supseteq T \times_n^+ T$ 



\end{frame}


\begin{frame}[t]{Surjective bounded morphism}
	Let $N_\omega^r \times_n^+ N_\omega^r = (X \times X, \tau_1, \tau_2, \tau)$ and $N(T_{\omega,\omega,\omega}^{r,m}) = ((\mathbb{N} \times \{0,1,2\})^*, \sigma_1, \sigma_2, \sigma)$.
	\begin{theorem}
	There is a surjective bounded morphism $g: N_\omega^r \times_n^+ N_\omega^r \to N(T_{\omega,\omega,\omega}^{r,m})$. 
	\end{theorem}
	Proof sketch: We fix a bijection $h: \mathbb{N} \times \mathbb{N} \to \mathbb{N} $. \pause
			\[
				g: X \times X \to (\mathbb{N} \times \{0,1,2\})^* 
		\] \pause
\vspace{-0.4cm}
			\[
\begin{array}{cccccc}
\alpha &= 3 & \star & 2 & \star & \star^\omega\\
\beta  &= \star     & 1& 7 & \star & \star^\omega\\
\hline \pause
g'(\alpha,\beta) &= (3,1) & (1,2) & (h(2,7),0) & \star & \star^\omega
\end{array}
\]
 \pause


Then $g(\alpha,\beta) = (3,1)(1,2) (h(2,7),0)$

	
\end{frame}





\begin{frame}{Surjective bounded morphism}
	Let $N_\omega^r \times_n^+ N_\omega^r = (X \times X, \tau_1, \tau_2, \tau)$ and $N(T_{\omega,\omega,\omega}^{r,m}) = ((\mathbb{N} \times \{0,1,2\})^*, \sigma_1, \sigma_2, \sigma)$ 
	\begin{theorem}
		
		There is a surjective bounded morphism $g: N_\omega^r \times_n^+ N_\omega^r \to N(T_{\omega,\omega,\omega}^{r,m})$.

		
	\end{theorem}
	%Pick $m > max\{st(\alpha),st(\beta)\}$.
	\begin{itemize}
		\item surjectivity and 
		\item $R(g(\alpha,\beta)) \subseteq g(U_k(\alpha) \times U_j(\beta)) \subseteq g(U)$ for some $k,j \in \mathbb{N}$
	\end{itemize}
		\vspace{1cm}
	    {\transparent{0.3}

   		2nd cond: for all $1 \le i \le k$, for all $w \in W$ and all $U \in \tau_i(w)$, there is a $V\in \sigma_i(g(w))$ with $V = g(U)$ 	
			Note: $\tau(x,y) = \{U \mid \exists k,j \in \mathbb{N}: U_k(x) \times U_j(y) \subseteq U\}, \sigma(x) = \{U \mid R(x) \subseteq U\}$.

			
    	}

\end{frame}

\begin{frame}{Surjective bounded morphism}


$g(\alpha,\beta) = (3,2)(1,0) h((7,2),0) \textcolor{red}{R} \mbox{     }(3,2)(1,0) h((7,2),0) \textcolor{red}{(h(6,7),0)}$

\[
\begin{array}{ccccccc}
\alpha' &= \star & 1 & 2 & \star & \textcolor{red}{6} & \star^\omega\\
\beta'  &= 3     & \star & 7 & \star & \textcolor{red}{7} & \star^\omega\\
\hline 
g'(\alpha,\beta) &= (3,2) & (1,0) & (h(2,7),0) & \star & \textcolor{red}{(h(6,7),0)}& \star^\omega
\end{array}
\]
\vspace{0.5cm}

{\transparent{0.3} 
	Claim: $R(g(\alpha,\beta)) \subseteq g(U_k(\alpha) \times U_j(\beta)) \subseteq g(U)$ for some $k,j \in \mathbb{N}$	}



	
\end{frame}

\begin{frame}{Surjective bounded morphism}
			Let $N_\omega^r \times_n^+ N_\omega^r = (X \times X, \tau_1, \tau_2, \tau)$ and $N(T_{\omega,\omega,\omega}^{r,m}) = ((\mathbb{N} \times \{0,1,2\} \cup \{\star\})^*, \sigma_1, \sigma_2, \sigma)$
		\begin{theorem}
		There is a surjective bounded morphism $g: N_\omega^r \times_n^+ N_\omega^r \to N(T_{\omega,\omega,\omega}^{r,m})$. Note: $Log(N_w^r) = T$ and hence $N_w^r \times_n^+ N_w^r \in T \times_n^+ T$
		\end{theorem}
		Pick $m > max\{st(\alpha),st(\beta)\}$ and set $U:= U_m(\alpha) \times U_m(\beta)$. Then 
		\[
			g(U) = g(U_m(\alpha) \times U_m(\beta)) \subseteq R(g(\alpha,\beta)) \subseteq V
		\]

		{\transparent{0.3}
		Note: $\tau(x,y) = \{U \mid \exists k,j \in \mathbb{N}: U_k(x) \times U_j(y) \subseteq U\}, \sigma(x) = \{U \mid R(x) \subseteq U\}$\newline
		3rd cond: for all $1 \le i \le k$, for all $w \in W$ and all $V \in \sigma_i(f(w))$, there exists $U \in \tau_i(w)$ such that $g(U) \subseteq V$ 
		}
\end{frame}

\begin{frame}[t]{Conlusion}
	
	
\begin{itemize}
	\setlength{\itemsep}{1em}

	\item $Log(T_{\omega,\omega,\omega}^{r,m}) = T_3m$ \pause
	\item $Log(T_{\omega,\omega,\omega}^{r,m}) \supseteq Log(N_\omega^w \times_n^+ N_\omega^w )\supseteq T \times_n^+ T $ \pause
	\item $T3m \supseteq T \times_n^+ T$
	
\end{itemize}



\end{frame}

\begin{frame}{Future works}
	\begin{itemize}
		\item $K \times_n^+ K, K4 \times_n^+K4, ...$	\pause	
		\item $T \times_n^+ D, K4 \times_n^+ D, S4 \times_n^+ K4...$
	\end{itemize}
\end{frame}

